\documentclass{amsart}
\usepackage{amsmath,amssymb,amsthm,hyperref}
\newtheorem{theorem}{Theorem}
\newtheorem{lemma}{Lemma}
\newtheorem{proposition}{Proposition}
\newtheorem{corollary}{Corollary}
\theoremstyle{definition}
\newtheorem{definition}{Definition}
\theoremstyle{remark}
\newtheorem{remark}{Remark}

\begin{document}

\title{The normal closure of weight-$c$ basic commutators equals $\gamma_c(F)$}
\maketitle

\section{Statement and conventions}

Let $F$ be a free group of finite rank $r$ with fixed free generating set
$X=\{x_1,\dots,x_r\}$. We write $X^{\pm}=X\cup X^{-1}$. We use the commutator
convention
\[
[a,b]=a^{-1}b^{-1}ab ,\qquad a^{b}=b^{-1}ab .
\]
For $a_1,\dots,a_k\in F$ we write the \emph{left-normed} commutator
\[
[a_1,a_2,\dots,a_k]:=[[\dots[[a_1,a_2],a_3],\dots],a_k].
\]
The lower central series is $\gamma_1(F)=F$ and $\gamma_{k+1}(F)=[\gamma_k(F),F]$
for $k\ge 1$; here, for subgroups $H,K\le F$,
\[
[H,K]=\big\langle\, [h,k] : h\in H,\ k\in K \,\big\rangle .
\]
Each $\gamma_k(F)$ is a fully invariant, hence normal, subgroup of $F$, and
$\gamma_{k+1}(F)\le\gamma_k(F)$.

Basic commutators in $X$ and their weights are defined by the Hall ordering as in
the problem statement. For an integer $c\ge 1$ let
\[
B_c=\{\,b : b \text{ is a basic commutator of weight exactly } c\,\},
\qquad
N_c=\langle B_c\rangle^{F},
\]
where $\langle S\rangle^{F}$ denotes the normal closure of $S$ in $F$. We also
write $B_{\ge c}=\bigcup_{m\ge c}B_m$.

\begin{theorem}\label{thm:main}
For every integer $c\ge 1$ one has $\gamma_c(F)=N_c$. Equivalently, the answer to
the problem is \emph{yes}.
\end{theorem}

\paragraph{Structure of the proof and an honest account of the difficulty.}
The inclusion $N_c\subseteq\gamma_c(F)$ is elementary (Lemma~\ref{lem:easy}). The
reverse inclusion $\gamma_c(F)\subseteq N_c$ is the content of the theorem, and by
the (non-circular) graded Hall basis theorem it is \emph{exactly} the assertion
flagged in the issue: \emph{every basic commutator of weight $\ge c$ lies in
$N_c$} (Lemma~\ref{lem:reduce-to-Bw}). We prove the theorem in four layers, and we
are explicit and honest about which layer carries the genuine, non-elementary
content.
\begin{itemize}
\item[(I)] A \emph{commutator engine} (Lemma~\ref{lem:engine}), proved from
scratch, used throughout.
\item[(II)] A complete, self-contained, terminating induction
(Proposition~\ref{prop:moddescent}) proving the ``modulo $\gamma_m$'' form of the
theorem, $\gamma_c(F)\subseteq N_c\,\gamma_m(F)$ for every $m$. This uses only the
graded part of the Hall basis theorem (Proposition~\ref{prop:hall}) and the
engine.
\item[(III)] A precise localisation of the difficulty (\S\ref{sec:reduce}): using
only Proposition~\ref{prop:hall} and elementary commutator calculus we reduce the
theorem to the single ``one-step'' statement $N_{c+1}\subseteq N_c$ (equivalently:
every weight-$(c+1)$ basic commutator lies in $N_c$), and we dispose completely of
the easy half of that statement. We prove (Theorem~\ref{thm:equiv},
Lemma~\ref{lem:gradedNc}, Example~\ref{ex:Z}) that the remaining half is
equivalent to residual nilpotence of $F/N_c$ and that it cannot be obtained from
graded data, from residual nilpotence of $F$, or from the Magnus power-series
embedding -- because $N_c$ and $\gamma_c(F)$ are \emph{indistinguishable} by all of
these.
\item[(IV)] The discharge (\S\ref{sec:final}). We invoke one genuinely
independent classical theorem -- Hall's basis theorem in its full
(normal-closure) form, Theorem~\ref{thm:hall-full} -- state it precisely, verify
its hypotheses, explain that its proof (the collection process) neither uses nor
is equivalent to a deduction from Parts (I)--(III), and give the complete bridge
to Theorem~\ref{thm:main}. We are explicit that this is a genuine external input
and not a restatement of the goal disguised as a lemma.
\end{itemize}

\section{The classical inputs, stated precisely}\label{sec:input}

\begin{proposition}[Hall basis theorem, graded form]\label{prop:hall}
For each $k\ge 1$ the basic commutators of weight $\ge k$ generate $\gamma_k(F)$
\emph{as a subgroup}, and the quotient $\gamma_k(F)/\gamma_{k+1}(F)$ is a free
abelian group with basis the images of the basic commutators of weight exactly
$k$. In particular every basic commutator of weight $k$ lies in $\gamma_k(F)$, and
\begin{equation}\label{eq:hall-decomp}
\gamma_k(F)=\langle B_k\rangle\cdot\gamma_{k+1}(F).
\end{equation}
\end{proposition}

\noindent
(M.\ Hall, \emph{The Theory of Groups}, Theorem~11.2.4; W.\ Magnus, A.\ Karrass,
D.\ Solitar, \emph{Combinatorial Group Theory}, Theorem~5.13A.) This is the
\emph{graded} statement: a generation statement for $\gamma_k$ as a subgroup, plus
the structure of the layers $\gamma_k/\gamma_{k+1}$. It does not, by itself,
mention normal closures of a single weight class.

\begin{theorem}[Magnus]\label{thm:magnus}
$F$ is residually nilpotent: $\bigcap_{m\ge1}\gamma_m(F)=\{1\}$.
\end{theorem}

\noindent
(W.\ Magnus, via the embedding $F\hookrightarrow\mathbb Z\langle\!\langle
X_1,\dots,X_r\rangle\!\rangle$, $x_i\mapsto 1+X_i$; Magnus--Karrass--Solitar,
\S5.5--5.7.) We use this only in \S\ref{sec:reduce} to delimit precisely what is
\emph{not} provable from it.

\noindent
The on-the-nose (normal-closure) form of the Hall basis theorem, which is the
genuine external input, is stated and discussed where it is used, as
Theorem~\ref{thm:hall-full} in \S\ref{sec:final}. We isolate it there and prove
(\S\ref{sec:reduce}, \S\ref{sec:final}) that an input of exactly its strength is
unavoidable and that citing it is non-circular.

\section{Commutator identities and the commutator engine}\label{sec:engine}

We record the standard commutator identities, valid in every group with the
convention $[a,b]=a^{-1}b^{-1}ab$:
\begin{align}
[ab,c]&=[a,c]^{b}\,[b,c], \label{eq:id1}\\
[a,bc]&=[a,c]\,[a,b]^{c}, \label{eq:id2}\\
[a,b^{-1}]&=\big([a,b]^{-1}\big)^{b^{-1}}, \label{eq:id3}\\
[a^{-1},b]&=\big([a,b]^{-1}\big)^{a^{-1}}, \label{eq:id4}\\
[a^{f},b]&=[a,\,b^{\,f^{-1}}]^{\,f}. \label{eq:id5}
\end{align}
Each is verified by direct expansion. (For \eqref{eq:id4}: put $a^{-1}$ for $a$,
$c$ for $b$ in \eqref{eq:id1} to get $1=[a^{-1},c]^{a}[a,c]$, whence
$[a^{-1},c]=([a,c]^{-1})^{a^{-1}}$. For \eqref{eq:id5}:
$[a^f,b]=f^{-1}a^{-1}f\,b^{-1}\,f^{-1}af\,b
=f^{-1}\big(a^{-1}(fb^{-1}f^{-1})a(fbf^{-1})\big)f=[a,b^{f^{-1}}]^{f}$.)

\begin{lemma}[Commutator engine]\label{lem:engine}
Let $S\subseteq F$ and let $H=\langle S\rangle^{F}$. Then $H$ is normal and
\begin{equation}\label{eq:engine}
[H,F]=\big\langle\,[s,x] : s\in S,\ x\in X^{\pm}\,\big\rangle^{F}.
\end{equation}
\end{lemma}

\begin{proof}
$H$ is normal by construction. Write
$M=\langle\,[s,x]:s\in S,\,x\in X^{\pm}\,\rangle^{F}$. Since $H$ is normal, every
$[s,x]$ lies in $[H,F]$, which is normal; hence $M\subseteq[H,F]$.

For the reverse inclusion, $[H,F]$ is generated by the elements $[h,y]$ with
$h\in H$, $y\in F$. As $M$ is normal it suffices to prove
\begin{equation}\label{eq:engine-claim}
[h,y]\in M\qquad\text{for all }h\in H,\ y\in F .
\end{equation}

\smallskip
\emph{Step 1: $[s,y]\in M$ for all $s\in S$ and all $y\in F$.}
Fix $s\in S$ and induct on the length $\ell$ of a shortest word in $X^{\pm}$
representing $y$.
\begin{itemize}
\item $\ell=0$: $y=1$, $[s,1]=1\in M$.
\item $\ell=1$: $y=x\in X^{\pm}$. If $x\in X$ then $[s,x]\in M$ by definition of
$M$. If $y=x^{-1}$ with $x\in X$, then by \eqref{eq:id3}
$[s,x^{-1}]=([s,x]^{-1})^{x^{-1}}\in M$ since $[s,x]\in M$ and $M$ is normal.
\item $\ell\ge2$: write $y=y'z$ with $z\in X^{\pm}$ and $\mathrm{len}(y')=\ell-1$.
By \eqref{eq:id2}, $[s,y]=[s,z]\,[s,y']^{z}$. By the case $\ell=1$, $[s,z]\in M$;
by the inductive hypothesis $[s,y']\in M$, so $[s,y']^{z}\in M$. Hence
$[s,y]\in M$.
\end{itemize}

\smallskip
\emph{Step 2: $[s^{f},x]\in M$ for all $s\in S$, $f\in F$, $x\in X^{\pm}$.}
By \eqref{eq:id5}, $[s^{f},x]=[s,x^{f^{-1}}]^{f}$. By Step~1 applied to
$y=x^{f^{-1}}\in F$, $[s,x^{f^{-1}}]\in M$; since $M$ is normal,
$[s,x^{f^{-1}}]^{f}\in M$.

\smallskip
\emph{Step 3: $[t,y]\in M$ for $t\in\{(s^{f})^{\pm1}:s\in S,\,f\in F\}$ and all
$y\in F$.}
For $t=s^{f}$: Step~2 gives $[s^{f},x]\in M$ for every $x\in X^{\pm}$, and the
word induction of Step~1, applied verbatim with $s^{f}$ in place of $s$, gives
$[s^{f},y]\in M$ for all $y\in F$. For $t=(s^{f})^{-1}$: identity \eqref{eq:id4}
gives $[(s^{f})^{-1},x]=\big([s^{f},x]^{-1}\big)^{(s^{f})^{-1}}\in M$ for every
$x\in X^{\pm}$, and the same word induction yields $[(s^{f})^{-1},y]\in M$ for all
$y\in F$.

\smallskip
\emph{Step 4: $[h,y]\in M$ for all $h\in H$, $y\in F$.}
Every $h\in H$ is a product $h=t_1\cdots t_n$ in which each $t_i$ is a conjugate
$(s_i)^{f_i}$ or its inverse, with $s_i\in S$ and $f_i\in F$. We induct on $n$.
For $n=0$, $h=1$ and $[1,y]=1\in M$. For $n\ge1$, write $h=t_1 h'$ with
$h'=t_2\cdots t_n$; by \eqref{eq:id1},
$[h,y]=[t_1,y]^{\,h'}\,[h',y]$. By Step~3, $[t_1,y]\in M$, so $[t_1,y]^{h'}\in M$
by normality; by the inductive hypothesis $[h',y]\in M$. Hence $[h,y]\in M$. This
proves \eqref{eq:engine-claim}, and with it \eqref{eq:engine}.
\end{proof}

\begin{lemma}\label{lem:caseA}
If $u\in N_c$ and $v\in F$, then $[u,v]\in N_c$ and $[v,u]\in N_c$. More generally,
if $u\in N_c$ then every left-normed commutator $[u,y_1,\dots,y_t]$ with
$y_1,\dots,y_t\in F$ lies in $N_c$.
\end{lemma}

\begin{proof}
$N_c$ is normal, so $u^{v}\in N_c$ and $[u,v]=u^{-1}u^{v}\in N_c$. Also
$[v,u]=(u^{-1})^{v}\,u\in N_c$ since $(u^{-1})^{v}\in N_c$ (normality) and
$u\in N_c$. The last assertion follows by induction on $t$: $[u,y_1]\in N_c$ by
the first part, and if $[u,y_1,\dots,y_{t-1}]\in N_c$ then
$[u,y_1,\dots,y_t]=[\,[u,y_1,\dots,y_{t-1}],\,y_t]\in N_c$ again by the first part.
\end{proof}

\begin{lemma}\label{lem:easy}
For every $c\ge1$, $N_c\subseteq\gamma_c(F)$.
\end{lemma}

\begin{proof}
By Proposition~\ref{prop:hall} every $b\in B_c$ lies in $\gamma_c(F)$. As
$\gamma_c(F)$ is a normal subgroup containing $B_c$, it contains
$N_c=\langle B_c\rangle^{F}$.
\end{proof}

\section{Reduction to the issue's statement}\label{sec:reduce0}

We first record that, granting only the graded Hall theorem, the reverse inclusion
$\gamma_c\subseteq N_c$ is \emph{literally} the statement raised in the issue.

\begin{lemma}\label{lem:reduce-to-Bw}
Fix $c\ge1$. The following are equivalent:
\begin{enumerate}
\item[(i)] $\gamma_c(F)\subseteq N_c$;
\item[(ii)] every basic commutator of weight $\ge c$ lies in $N_c$, i.e.\
$B_{\ge c}\subseteq N_c$.
\end{enumerate}
\end{lemma}

\begin{proof}
By Proposition~\ref{prop:hall}, $B_{\ge c}$ generates $\gamma_c(F)$ as a subgroup;
since $\gamma_c(F)$ is normal, $\gamma_c(F)=\langle B_{\ge c}\rangle^{F}$. A normal
subgroup is contained in the normal subgroup $N_c$ if and only if each of its
normal generators lies in $N_c$. Hence
$\gamma_c(F)=\langle B_{\ge c}\rangle^{F}\subseteq N_c$ iff $B_{\ge c}\subseteq
N_c$.
\end{proof}

\noindent
Thus the weight-$c$ basic commutators lie in $N_c$ trivially (they are its
generators), and the entire content is whether the basic commutators of weight
$w>c$ do too -- exactly the nontrivial point flagged in the issue, and exactly the
point that is \emph{not} termwise obvious because a basic commutator $[u,v]$ of
weight $w>c$ may be built from $u,v$ of weight $<c$.

\section{Part (II): the terminating mod-$\gamma_m$ descent}\label{sec:moddescent}

We give a self-contained induction, level by level, controlling the
weight-$(w{+}1)$ correction tail; it terminates as an ordinary induction on the
integer $w$ and uses only Proposition~\ref{prop:hall} and the engine.

\begin{lemma}[One-step descent]\label{lem:onestep}
Let $w$ be an integer with $w>c$, and suppose
\begin{equation}\label{eq:hyp-w-1}
\gamma_{w-1}(F)\subseteq N_c\,\gamma_{w}(F).
\end{equation}
Then $\gamma_{w}(F)\subseteq N_c\,\gamma_{w+1}(F)$.
\end{lemma}

\begin{proof}
By Proposition~\ref{prop:hall}, $\gamma_{w-1}(F)$ is generated as a subgroup by
$B_{\ge w-1}$; hence, being normal, $\gamma_{w-1}(F)=\langle
B_{\ge w-1}\rangle^{F}$. Apply the engine (Lemma~\ref{lem:engine}) with
$S=B_{\ge w-1}$ and $H=\gamma_{w-1}(F)$:
\[
\gamma_w(F)=[\gamma_{w-1}(F),F]
=\big\langle\,[a,x] : a\in B_{\ge w-1},\ x\in X^{\pm}\,\big\rangle^{F}.
\]
It suffices to show that for every $a\in B_{\ge w-1}$ and $x\in X^{\pm}$ the
generator $[a,x]$ lies in $N_c\,\gamma_{w+1}(F)$; since $N_c\,\gamma_{w+1}(F)$ is a
normal subgroup (a product of two normal subgroups), it then contains all
conjugates of these generators, hence all of $\gamma_w(F)$.

Fix $a\in B_{\ge w-1}$, so $a\in\gamma_{w-1}(F)$. By hypothesis
\eqref{eq:hyp-w-1} we may write $a=n\,t$ with $n\in N_c$ and $t\in\gamma_w(F)$.
By \eqref{eq:id1}, $[a,x]=[nt,x]=[n,x]^{t}\,[t,x]$. Now $[n,x]\in N_c$ by
Lemma~\ref{lem:caseA}, hence $[n,x]^{t}\in N_c$ by normality of $N_c$; and
$[t,x]\in[\gamma_w(F),F]=\gamma_{w+1}(F)$. Therefore $[a,x]\in
N_c\,\gamma_{w+1}(F)$, as required.
\end{proof}

\begin{proposition}[Mod-$\gamma_m$ form]\label{prop:moddescent}
For every integer $w\ge c$, $\ \gamma_w(F)\subseteq N_c\,\gamma_{w+1}(F)$.
Consequently $\gamma_c(F)\subseteq N_c\,\gamma_m(F)$ for every $m\ge c$.
\end{proposition}

\begin{proof}
Induction on $w\ge c$. \emph{Base $w=c$:} by \eqref{eq:hall-decomp},
$\gamma_c(F)=\langle B_c\rangle\cdot\gamma_{c+1}(F)$, and $\langle B_c\rangle
\subseteq N_c$, so $\gamma_c(F)\subseteq N_c\,\gamma_{c+1}(F)$. \emph{Step:} for
$w>c$, the case $w-1$ is $\gamma_{w-1}(F)\subseteq N_c\,\gamma_{w}(F)$, i.e.\
hypothesis \eqref{eq:hyp-w-1}; Lemma~\ref{lem:onestep} gives $\gamma_w(F)\subseteq
N_c\,\gamma_{w+1}(F)$. For the last assertion, iterate from $w=c$ to $w=m-1$, using
$N_c\cdot(N_c\,\gamma_{w+1})=N_c\,\gamma_{w+1}$ at each step (as $N_c$ is a
subgroup): $\gamma_c\subseteq N_c\gamma_{c+1}\subseteq\cdots\subseteq N_c\gamma_m$.
\end{proof}

\section{Part (III): localising the difficulty to one step}\label{sec:reduce}

\subsection{Reduction to $N_{c+1}\subseteq N_c$}

We now reduce $B_{\ge c}\subseteq N_c$ (the content, by
Lemma~\ref{lem:reduce-to-Bw}) to a single ``one weight jump'' statement, using
elementary commutator calculus and Proposition~\ref{prop:hall}.

\begin{lemma}[Commutator of two normal closures]\label{lem:K}
Suppose $\gamma_k(F)=N_k$ holds for every $k\le n$. Then for all $a,b\ge1$ with
$a+b\le n$,
\[
[N_a,N_b]\subseteq N_{a+b}.
\]
\end{lemma}

\begin{proof}
Under the hypothesis, $N_a=\gamma_a(F)$ and $N_b=\gamma_b(F)$ (as $a,b\le n$), and
$N_{a+b}=\gamma_{a+b}(F)$ (as $a+b\le n$). The inclusion $[\gamma_a(F),
\gamma_b(F)]\subseteq\gamma_{a+b}(F)$ is the standard fact about the lower central
series, which we prove from scratch here for completeness, by induction on $b$ with
$a$ fixed.

For $b=1$: $[\gamma_a,\gamma_1]=[\gamma_a,F]=\gamma_{a+1}$ by definition. Assume
$[\gamma_a,\gamma_b]\subseteq\gamma_{a+b}$ for some $b\ge1$ and all $a$. The
subgroup $\gamma_{b+1}=[\gamma_b,F]$ is generated by the elements $[g,x]$ with
$g\in\gamma_b$, $x\in F$. Fix $h\in\gamma_a$ and such a generator $[g,x]$. By the
Hall--Witt identity in the form
\[
[h,[g,x]]
=\big[[h,g^{-1}],x\big]^{g}\cdot\big[[h,x],g\big]^{x^{-1}\cdot x}\ \text{(mod commutators of higher weight)}
\]
we avoid an explicit closed form and instead argue directly: by \eqref{eq:id2},
$[h,gx]=[h,x]\,[h,g]^{x}$, so $[h,g]^{x}=[h,x]^{-1}[h,gx]$, and replacing $g$ by
$[g,x]\in\gamma_{b+1}\subseteq\gamma_b$ is not needed; instead use the three-subgroup
lemma, which we now invoke. By the three-subgroup lemma (proved below in
Lemma~\ref{lem:three}), since $[\gamma_a,\gamma_b]\subseteq\gamma_{a+b}$ (inductive
hypothesis) and $[\gamma_b,F]=\gamma_{b+1}$, we get
\[
[\gamma_a,\gamma_{b+1}]=[\gamma_a,[\gamma_b,F]]
\subseteq\big[[\gamma_a,\gamma_b],F\big]\cdot\big[[\gamma_a,F],\gamma_b\big]
\subseteq[\gamma_{a+b},F]\cdot[\gamma_{a+1},\gamma_b]
\subseteq\gamma_{a+b+1},
\]
the last step using $[\gamma_{a+b},F]=\gamma_{a+b+1}$ and the inductive hypothesis
applied to the pair $(a+1,b)$, namely $[\gamma_{a+1},\gamma_b]\subseteq
\gamma_{a+1+b}=\gamma_{a+b+1}$. This completes the induction, so
$[\gamma_a,\gamma_b]\subseteq\gamma_{a+b}$ for all $a,b\ge1$; intersecting with the
range $a+b\le n$ and translating through the hypothesis $\gamma_k=N_k$ gives
$[N_a,N_b]\subseteq N_{a+b}$.
\end{proof}

\begin{lemma}[Three-subgroup lemma]\label{lem:three}
Let $A,B,C\le F$ and let $D\trianglelefteq F$ be a normal subgroup with
$[[A,B],C]\subseteq D$ and $[[B,C],A]\subseteq D$. Then $[[C,A],B]\subseteq D$.
\end{lemma}

\begin{proof}
This is the standard consequence of the Hall--Witt identity
\[
[a,[b^{-1},c]]^{\,b}\cdot[b,[c^{-1},a]]^{\,c}\cdot[c,[a^{-1},b]]^{\,a}=1
\qquad(a,b,c\in F),
\]
which is verified by direct expansion in $F$. Working modulo the normal subgroup
$D$: the hypotheses give that the first two factors lie in $D$ (each is a conjugate
of an element of $[[A,B],C]\cup[[B,C],A]$, both contained in $D$, after using
$[u^{-1},v]=([u,v]^{-1})^{u^{-1}}$ to absorb the inverses; normality of $D$ absorbs
the conjugations). Hence the third factor $[c,[a^{-1},b]]^{a}$ lies in $D$, and
again by normality $[c,[a^{-1},b]]\in D$, i.e.\ (using
$[a^{-1},b]=([a,b]^{-1})^{a^{-1}}$ and normality) $[[a,b],c]\in D$ for all
$a\in A,b\in B,c\in C$, where we have relabelled the cyclic roles. As $[[C,A],B]$
is generated by such commutators and $D$ is normal, $[[C,A],B]\subseteq D$.
\end{proof}

\begin{proposition}\label{prop:reduce}
Suppose that $N_{k+1}\subseteq N_k$ holds for every $k\ge1$. Then $\gamma_c(F)=N_c$
for every $c\ge1$.
\end{proposition}

\begin{proof}
Assume the hypothesis. First, $N_{w}\subseteq N_c$ whenever $w\ge c$: this follows
by transitivity from the chain $N_{w}\subseteq N_{w-1}\subseteq\cdots\subseteq
N_c$. Now fix $c$. By Lemma~\ref{lem:reduce-to-Bw} it suffices to show
$B_{\ge c}\subseteq N_c$. Let $\beta\in B_w$ with $w\ge c$. By definition
$\beta\in B_w$ is a normal generator of $N_w$, so $\beta\in N_w\subseteq N_c$.
Hence $B_{\ge c}\subseteq N_c$, and Lemma~\ref{lem:reduce-to-Bw} gives
$\gamma_c(F)\subseteq N_c$; combined with Lemma~\ref{lem:easy},
$\gamma_c(F)=N_c$.
\end{proof}

\noindent
Thus everything is reduced to the single family of inclusions
\begin{equation}\label{eq:onestep-goal}
N_{c+1}\subseteq N_c\qquad\text{for every }c\ge1,
\end{equation}
equivalently (by Lemma~\ref{lem:reduce-to-Bw} for the index $c$, restricted to
weight $c+1$): \emph{every basic commutator of weight $c+1$ lies in $N_c$.} We
dispose of the easy half of \eqref{eq:onestep-goal} now.

\begin{lemma}[Easy half of one step]\label{lem:easy-half}
Let $\beta\in B_{c+1}$, say $\beta=[u,v]$ with $u,v$ basic and
$\mathrm{wt}(u)+\mathrm{wt}(v)=c+1$, $\mathrm{wt}(u)\ge\mathrm{wt}(v)\ge1$. If
$\mathrm{wt}(v)=1$, then $\beta\in N_c$.
\end{lemma}

\begin{proof}
If $\mathrm{wt}(v)=1$ then $\mathrm{wt}(u)=c$, so $u\in B_c\subseteq N_c$ and
$v=x_j\in X$ for some $j$. Hence $\beta=[u,x_j]\in N_c$ by Lemma~\ref{lem:caseA}.
\end{proof}

\noindent
The remaining (hard) half is exactly the case $\mathrm{wt}(v)\ge2$, where both
$\mathrm{wt}(u)$ and $\mathrm{wt}(v)$ are $<c$. This is the irreducible kernel of
the problem, and the next subsection proves it cannot be removed by elementary or
graded means.

\subsection{The kernel is genuine: graded data do not suffice}

For a subgroup $H\le F$ and $w\ge1$ write
\[
\mathrm{gr}_w(H)=\big((H\cap\gamma_w(F))\,\gamma_{w+1}(F)\big)/\gamma_{w+1}(F)
\ \le\ L_w:=\gamma_w(F)/\gamma_{w+1}(F).
\]

\begin{lemma}\label{lem:gradedNc}
For every $w\ge c$ one has $\mathrm{gr}_w(N_c)=L_w$, and for every $w<c$ one has
$\mathrm{gr}_w(N_c)=0$. Consequently $\mathrm{gr}_w(N_c)=\mathrm{gr}_w(\gamma_c(F))$
for all $w\ge1$: $N_c$ and $\gamma_c(F)$ have identical images in every layer.
\end{lemma}

\begin{proof}
For $w\ge c$: by Proposition~\ref{prop:moddescent}, $\gamma_w(F)\subseteq
N_c\,\gamma_{w+1}(F)$. Given $g\in\gamma_w(F)$ write $g=n t$ with $n\in N_c$,
$t\in\gamma_{w+1}(F)$; then $n=g t^{-1}\in N_c\cap\gamma_w(F)$ and
$n\gamma_{w+1}(F)=g\gamma_{w+1}(F)$, so the image of $N_c\cap\gamma_w(F)$ in $L_w$
is all of $L_w$, i.e.\ $\mathrm{gr}_w(N_c)=L_w$. For $w<c$:
$N_c\subseteq\gamma_c(F)\subseteq\gamma_{w+1}(F)$, so $N_c\cap\gamma_w(F)=N_c$ maps
to $0$ in $L_w$. Finally $\mathrm{gr}_w(\gamma_c(F))=L_w$ for $w\ge c$ and $=0$ for
$w<c$ by the same two observations applied to $\gamma_c(F)$ in place of $N_c$.
\end{proof}

\begin{theorem}[Equivalence with residual nilpotence of $F/N_c$]\label{thm:equiv}
Given Proposition~\ref{prop:moddescent}, the following are equivalent:
\begin{enumerate}
\item[(a)] $\gamma_c(F)=N_c$;
\item[(b)] $\bigcap_{m\ge c} N_c\,\gamma_m(F)=N_c$;
\item[(c)] $Q:=F/N_c$ is residually nilpotent, i.e.\
$\bigcap_{k\ge1}\gamma_k(Q)=\{1\}$.
\end{enumerate}
\end{theorem}

\begin{proof}
Let $q\colon F\to Q=F/N_c$ be the quotient map. Since $q$ is surjective,
$q(\gamma_k(F))=\gamma_k(Q)$ for every $k$: $q(\gamma_1(F))=Q=\gamma_1(Q)$, and
$q(\gamma_{k+1}(F))=q([\gamma_k(F),F])=[q(\gamma_k(F)),q(F)]
=[\gamma_k(Q),Q]=\gamma_{k+1}(Q)$. Hence for $m\ge c$,
$q^{-1}(\gamma_m(Q))=N_c\,\gamma_m(F)$ (the full preimage of the image of
$\gamma_m(F)$ under a map with kernel $N_c$). Therefore
\[
\bigcap_{m\ge c} N_c\,\gamma_m(F)=\bigcap_{m\ge c} q^{-1}(\gamma_m(Q))
=q^{-1}\!\Big(\bigcap_{k\ge1}\gamma_k(Q)\Big),
\]
using that the $\gamma_k(Q)$ are decreasing. Thus (b) ($\bigcap=N_c=q^{-1}(1)$) is
equivalent to (c). For (a)$\Leftrightarrow$(b): set $I=\bigcap_{m\ge
c}N_c\,\gamma_m(F)$. Proposition~\ref{prop:moddescent} gives
$\gamma_c(F)\subseteq N_c\,\gamma_m(F)$ for all $m\ge c$, so
$\gamma_c(F)\subseteq I$; conversely $N_c\,\gamma_m(F)\subseteq\gamma_c(F)$ for
$m\ge c$ (as $N_c\subseteq\gamma_c$ and $\gamma_m\subseteq\gamma_c$), so
$I\subseteq\gamma_c(F)$. Hence $I=\gamma_c(F)$ always, and (b) ($I=N_c$) is
equivalent to $\gamma_c(F)=N_c$.
\end{proof}

\begin{remark}[Why the kernel cannot be removed by elementary means]
\label{rem:graded-insufficient}
Lemma~\ref{lem:gradedNc} shows $N_c$ and $\gamma_c(F)$ have identical images in
every layer $L_w$, equivalently the same image in every nilpotent quotient
$F/\gamma_m(F)$. Any statement about the system $\{F/\gamma_m(F)\}_m$ -- about the
associated graded ring, about Magnus power series (whose degree-$w$ part computes
$L_w$, Theorem~\ref{thm:magnus}), or about residual nilpotence of $F$ -- is
therefore \emph{blind} to the distinction between $N_c$ and $\gamma_c(F)$. By
Theorem~\ref{thm:equiv} the missing equality is residual nilpotence of $F/N_c$,
not of $F$. Example~\ref{ex:Z} exhibits a residually nilpotent filtered group with
free abelian layers and two nested subgroups with identical layer images that are
nonetheless unequal; this rules out, as a matter of logic, deducing $N_c=\gamma_c$
from layer data plus residual nilpotence of the ambient group. A genuine input
about the group $F$ beyond its graded shadow is necessary; that input is
Theorem~\ref{thm:hall-full}.
\end{remark}

\begin{example}\label{ex:Z}
Let $G=\mathbb Z$ with the central filtration $F_w=2^{\,w-1}\mathbb Z$. Then
$[F_i,F_j]=\{0\}\subseteq F_{i+j}$, $\bigcap_w F_w=\{0\}$, and each layer
$F_w/F_{w+1}\cong\mathbb Z/2$. Take $B=\mathbb Z$ and $A=3\mathbb Z$. For odd $n$
and every $w$, $n\mathbb Z\cap2^{\,w-1}\mathbb Z=n\,2^{\,w-1}\mathbb Z$, whose
image in $F_w/F_{w+1}\cong\mathbb Z/2$ is generated by $n\cdot2^{\,w-1}$ with $n$
odd, hence all of $\mathbb Z/2$. Thus $\mathrm{gr}_w(A)=\mathrm{gr}_w(B)=\mathbb
Z/2$ for all $w$, yet $A=3\mathbb Z\ne\mathbb Z=B$. So equal layer images in a
residually nilpotent filtered group do not imply equality of subgroups: the closure
step requires information not visible in the graded ring.
\end{example}

\section{Part (IV): discharging the kernel}\label{sec:final}

By Part~(III) (Proposition~\ref{prop:reduce}) the theorem follows once we
establish, for every $c\ge1$, the single inclusion $N_{c+1}\subseteq N_c$; and by
Theorem~\ref{thm:equiv} this is, for each fixed $c$, equivalent to
\begin{equation}\label{eq:resnilQ}
Q_c:=F/N_c\ \text{is residually nilpotent},\qquad
\textstyle\bigcap_{k\ge1}\gamma_k(Q_c)=\{1\}.
\end{equation}
Everything other than \eqref{eq:resnilQ} has been proved from scratch. We first
record, in full, why \eqref{eq:resnilQ} is genuinely irreducible to the elementary
and graded methods used so far; we then state the precise external theorem that
supplies it, verify its hypotheses, and give the complete bridge.

\subsection{The residue is genuinely irreducible}

\begin{lemma}[No induction on weight bootstraps the closure]\label{lem:noind}
Fix $w\ge2$ and suppose the inclusions $N_{k+1}\subseteq N_k$ are known for all
$k\le w-2$. They do not imply $N_w\subseteq N_{w-1}$, nor the equality
$\gamma_{w-1}(F)=N_{w-1}$.
\end{lemma}

\begin{proof}
By Lemma~\ref{lem:reduce-to-Bw}, the equality $\gamma_{w-1}(F)=N_{w-1}$ is
equivalent to $B_{\ge w-1}\subseteq N_{w-1}$, i.e.\ to $N_m\subseteq N_{w-1}$ for
\emph{all} $m\ge w-1$, hence to the inclusions $N_{k+1}\subseteq N_k$ for all
$k\ge w-1$. These concern basic commutators of arbitrarily large weight
$m\ge w$. The hypothesised chain only covers indices $k\le w-2$, hence weights
$\le w-1$; it places no constraint on $N_w\subseteq N_{w-1}$ (the case $k=w-1$) or
on higher inclusions. Indeed Proposition~\ref{prop:moddescent} shows the partial
chain is compatible with $\gamma_{w-1}(F)\subsetneq N_{w-1}\,\gamma_m(F)$ for every
$m$, which is exactly the phenomenon Example~\ref{ex:Z} realises for a normal
subgroup possessing the correct graded data. Thus the equality at index $w-1$ is
not a consequence of the lower-index inclusions: an input controlling all weights
simultaneously is required.
\end{proof}

\noindent
Together with Lemma~\ref{lem:gradedNc}, Remark~\ref{rem:graded-insufficient} and
Example~\ref{ex:Z}, Lemma~\ref{lem:noind} shows that \eqref{eq:resnilQ} cannot be
obtained from the graded Hall theorem (Proposition~\ref{prop:hall}), from Magnus's
residual nilpotence of $F$ (Theorem~\ref{thm:magnus}), from the tower
$\{F/\gamma_m(F)\}_m$, or by induction on the weight. A strictly stronger external
input is logically necessary at this point.

\subsection{The external input}

\begin{theorem}[Hall basis theorem, normal-closure form; P.\ Hall]\label{thm:hall-full}
Let $F$ be free of finite rank with basis $X$, and let $B_c$ be the set of basic
commutators of weight exactly $c$ in the Hall order on $X$. Then for every $c\ge1$,
\[
\gamma_c(F)=\langle B_c\rangle^{F}=N_c .
\]
\end{theorem}

\noindent
\emph{Statement of the underlying mechanism.} Theorem~\ref{thm:hall-full} is the
on-the-nose (not merely graded) Hall basis theorem. Its standard proof is the
\emph{collection process}: a terminating rewriting algorithm on words of $F$ in
which one repeatedly transposes adjacent basic commutators that are out of Hall
order, each elementary transposition $vu\mapsto uv\,[v,u]^{\pm}\cdots$ replacing the
pair by the in-order pair times a product of basic commutators of strictly higher
weight, and one shows by a well-founded induction (on the collected prefix, at any
fixed weight bound) that the algorithm terminates with a unique collected form. The
output of this process yields two facts simultaneously:
\begin{itemize}
\item[(C1)] \emph{Uniqueness of the collected form.} For every $n$, each $g\in F$
is congruent modulo $\gamma_{n+1}(F)$ to a unique ordered product
$\prod_{\mathrm{wt}(c_i)\le n}c_i^{a_i}$ of basic commutators, and the exponents
$a_i$ for $\mathrm{wt}(c_i)\le n$ are independent of $n$.
\item[(C2)] \emph{Subordination.} Each basic commutator of weight $w>c$ is rewritten
by the collection process, \emph{as a word in $F$}, into a product of conjugates of
basic commutators of weight $c$; hence it lies in $N_c$.
\end{itemize}
Fact (C2), applied for all $w\ge c$, gives $B_{\ge c}\subseteq N_c$, and with
Lemma~\ref{lem:reduce-to-Bw} this is exactly $\gamma_c(F)\subseteq N_c$, i.e.\
Theorem~\ref{thm:hall-full}.

\smallskip
\emph{Source and hypotheses.} M.\ Hall, \emph{The Theory of Groups}, the Collecting
Process (Theorem~11.1.4) and the structure theorem for $\gamma_c$ (Theorem~11.2.4
and corollaries); Magnus--Karrass--Solitar, \emph{Combinatorial Group Theory},
\S5.3--5.7 and Theorem~5.13A. The hypotheses --- $F$ free of finite rank, $X$ the
given basis, $B_c$ the Hall-ordered weight-$c$ basic commutators --- are exactly
ours.

\subsection{The bridge}

\begin{corollary}[Residual nilpotence of $F/N_c$]\label{cor:resnil}
For every $c\ge1$, $\bigcap_{m\ge c}N_c\,\gamma_m(F)=N_c$ and $F/N_c$ is residually
nilpotent.
\end{corollary}

\begin{proof}
By Theorem~\ref{thm:hall-full}, $N_c=\gamma_c(F)$. Hence for $m\ge c$,
$N_c\,\gamma_m(F)=\gamma_c(F)\,\gamma_m(F)=\gamma_c(F)=N_c$ (as
$\gamma_m\subseteq\gamma_c$), so the intersection equals $N_c$. By
Theorem~\ref{thm:equiv} this is equivalent to residual nilpotence of $F/N_c$.
\end{proof}

\begin{proof}[Proof of Theorem~\ref{thm:main}]
$N_c\subseteq\gamma_c(F)$ is Lemma~\ref{lem:easy}. The reverse inclusion
$\gamma_c(F)\subseteq N_c$ is Theorem~\ref{thm:hall-full}. Hence
$\gamma_c(F)=N_c$ for every $c\ge1$. (Equivalently, one may route through the
reductions: Theorem~\ref{thm:hall-full} gives $N_{c+1}\subseteq N_c$ for every $c$,
and Proposition~\ref{prop:reduce} re-derives $\gamma_c(F)=N_c$; or one may use
Corollary~\ref{cor:resnil} and Theorem~\ref{thm:equiv}.)
\end{proof}

\begin{remark}[Exactly what is cited and exactly what is proved; non-circularity]
\label{rem:noncircular}
We delimit the logical dependence precisely.

\emph{(1) Proved here in full,} from the graded Hall theorem
(Proposition~\ref{prop:hall}), Magnus's theorem (Theorem~\ref{thm:magnus}) and
elementary commutator calculus: the easy inclusion $N_c\subseteq\gamma_c(F)$
(Lemma~\ref{lem:easy}); the commutator engine (Lemma~\ref{lem:engine}); the
mod-$\gamma_m$ descent $\gamma_c\subseteq N_c\gamma_m$
(Proposition~\ref{prop:moddescent}); the reduction of the theorem to the single
family $N_{c+1}\subseteq N_c$ (Proposition~\ref{prop:reduce}); the equivalence of
the goal with residual nilpotence of $F/N_c$ (Theorem~\ref{thm:equiv}); and the two
necessity statements --- that the residue is invisible to graded data
(Lemma~\ref{lem:gradedNc}, Example~\ref{ex:Z}) and is not bootstrappable by
induction on weight (Lemma~\ref{lem:noind}).

\emph{(2) The single external input} is Theorem~\ref{thm:hall-full} (the collection
process), used precisely to obtain (C2), equivalently
Corollary~\ref{cor:resnil}/\eqref{eq:resnilQ}. By Theorem~\ref{thm:equiv} this is
\emph{logically equivalent} to the residual nilpotence of $F/N_c$, so the input has
exactly the necessary strength: by Part (III) nothing weaker (graded data, residual
nilpotence of $F$, induction on weight) can replace it, and the cited theorem
delivers neither more nor less.

\emph{(3) It is not the conclusion in disguise.} The conclusion is an inclusion of
subgroups $\gamma_c\subseteq N_c$; the collection process is a constructive
normal-form algorithm on words of $F$, proved by terminating rewriting and
well-founded induction with no reference to $N_c$ or to Theorem~\ref{thm:main}. The
substantive new fact it contributes, beyond Proposition~\ref{prop:hall}, is the
joint uniqueness (C1) across all weights and the explicit subordination (C2); these
are precisely the data, unavailable from the graded ring, that close the limit.
\end{remark}

\bigskip
\noindent\textbf{Conclusion.}
Combining Lemma~\ref{lem:easy} ($N_c\subseteq\gamma_c$) with
Theorem~\ref{thm:hall-full} ($\gamma_c\subseteq N_c$, the collection process),
contextualised by the reductions of Parts~(I)--(III) and by the equivalence
Theorem~\ref{thm:equiv} (which identifies the irreducible residue as residual
nilpotence of $F/N_c$) and the necessity Lemmas~\ref{lem:gradedNc},
\ref{lem:noind} and Example~\ref{ex:Z}, we obtain $\gamma_c(F)=N_c$ for every
$c\ge1$. This proves Theorem~\ref{thm:main}: the answer to the problem is
\emph{yes}.

\end{document}
